\section{Introduction}
\label{sec:introduction}

Consider the following scenario: a developer asks an SWE agent to ``fix this bug in the login flow.'' The agent begins investigating, locates a suspicious null-pointer dereference in the session handler, and proposes a patch. The developer reviews the diff and clarifies: ``actually, the issue is in the auth module---the token validation is wrong.'' The agent pivots, now examining authentication logic. Mid-way through the fix, the developer adds a constraint: ``don't change the public API---we have downstream consumers.'' The agent revises its approach accordingly. After seeing the updated patch, the developer reconsiders: ``use the other approach you mentioned earlier---the one with the middleware wrapper.'' The agent must now backtrack, discard its current changes, and reconstruct a solution that honors all accumulated constraints.

In existing SWE benchmarks~\citep{jimenez2024swebench, yang2024sweagent}, this interaction would never arise. SWE-bench evaluates agents on isolated GitHub issues: one issue description maps to one gold patch, with no user interaction whatsoever. This single-turn, static-intent formulation has driven remarkable progress---yet it systematically exempts agents from one of the most fundamental challenges of real-world software engineering assistance: keeping up with a developer whose understanding of the problem evolves as the agent works on it. We term this phenomenon \textbf{\drift{}}---the dynamic evolution of user intent during multi-turn interactions with SWE agents.

\paragraph{The static intent assumption in SWE benchmarks.}
The dominant evaluation paradigm for SWE agents assumes that the developer's intent is fully and correctly specified at the outset. SWE-bench~\citep{jimenez2024swebench} provides a single issue description and evaluates the agent's patch against repository test suites, with no opportunity for clarification or refinement. SWE-agent~\citep{yang2024sweagent} introduces a sophisticated agent-computer interface but inherits the same single-turn evaluation protocol. More recent efforts have begun to relax certain static assumptions: SWE-EVO~\citep{sweevo2025} studies cross-version evolution of codebases, testing whether agents can adapt patches as repositories change over time, but does not model interactive user intent changes within a single task. SWE-CI~\citep{sweci2026} introduces continuous integration feedback loops, but the requirements are generated from automated CI signals rather than from an evolving human developer. Neither benchmark captures the iterative, conversational process through which developers and SWE agents collaboratively arrive at solutions.

\paragraph{Why \drift{} matters for SWE agents.}
Intent drift poses uniquely severe challenges in software engineering contexts. Unlike conversational assistants that produce text responses, SWE agents execute \emph{concrete, often irreversible actions}: they edit source files, delete code, refactor APIs, modify configurations, and run test suites whose side effects persist. When a developer's intent shifts after the agent has already committed changes---renaming functions that other modules depend on, deleting files that turn out to be needed, or restructuring code in ways that are difficult to undo---the agent faces a fundamentally harder problem than simply updating its next response. It must reason about which prior edits remain valid under the new intent, which must be rolled back, and how to reconcile the revised goal with the current repository state. A failure to handle such drift gracefully can leave the codebase in an inconsistent state, introduce subtle regressions, or silently violate constraints the developer assumed were being respected. These failure modes are invisible to benchmarks that assume static, single-turn intents.

\paragraph{Our contributions.}
We introduce \textbf{\drift{}}, a systematic framework for evaluating SWE agents under dynamic user intent evolution. Our work makes the following contributions:

\begin{itemize}[leftmargin=*, itemsep=2pt]
    \item \textbf{A taxonomy of intent drift} (\S\ref{sec:taxonomy}): We identify and formalize six distinct patterns of intent drift---progressive refinement, topic shift, constraint addition, goal conflict, implicit need emergence, and intent backtracking---grounded in HCI literature and analysis of real developer--agent interactions. Each drift event is further decomposed along five orthogonal dimensions (distance, explicitness, timing, frequency, reversibility), enabling fine-grained characterization.

    \item \textbf{\bench{}-SWE} (\S\ref{sec:benchmark}): We construct a benchmark of 120 multi-turn SWE tasks spanning 6 popular Python repositories, each annotated with drift types, drift points, per-turn intent states, and expected repository end-states verified by test suites. Tasks are constructed via a two-stage \emph{blueprint-to-trajectory} pipeline with three-layer verification (automated test validation, cross-model LLM review, and dual human annotation with Cohen's $\kappa \geq 0.75$).

    \item \textbf{Deterministic, drift-aware metrics} (\S\ref{sec:metrics}): We propose five evaluation metrics grounded in test-based, deterministic assessment: \emph{pass\textsuperscript{k}} (adapted from $\tau$-bench~\citep{yao2024taubench}), measuring the probability that all $k$ sampled trajectories pass the final test suite; \emph{Drift Detection Latency} (DDL), quantifying turns elapsed before the agent acknowledges a drift event; \emph{Context Preservation Score} (CPS), assessing retention of still-valid prior context across drift boundaries; \emph{Rollback Rate}, measuring how often the agent correctly undoes invalidated changes; and \emph{Proactive Clarification Rate} (PCR), capturing unprompted requests for disambiguation.

    \item \textbf{Systematic evaluation and findings} (\S\ref{sec:experiments}): We evaluate multiple frontier models and SWE agent frameworks, including SWE-agent~\citep{yang2024sweagent}, Claude Code, and Cursor. Key findings include: (1) intent drift causes substantial pass\textsuperscript{k} degradation compared to static-intent baselines, (2) constraint addition and intent backtracking are the most challenging drift types, as they require reasoning about action reversibility in code, (3) implicit drift is significantly harder than explicit drift across all agents, and (4) agents with richer context management (\eg, file-level undo, checkpoint mechanisms) show disproportionate advantages on backtracking scenarios.
\end{itemize}

\paragraph{Positioning.}
We distinguish \drift{} from adjacent efforts along two axes. First, unlike SWE-EVO~\citep{sweevo2025} and SWE-CI~\citep{sweci2026}, which study agent robustness to \emph{codebase-level} or \emph{CI-driven} changes, we focus on \emph{user-side} intent evolution within a single interactive session. Second, unlike $\tau$-bench~\citep{yao2024taubench} and $\tau$-Knowledge~\citep{tauknowledge2026}, which evaluate tool agents in customer-service domains with simulated users, we target the software engineering domain where actions are persistent, state-modifying, and subject to hard correctness constraints via test suites. Our benchmark and metrics are designed to expose the gap between static evaluations and the interactive reality of SWE agent deployment.
